% Part: normal-modal-logic
% Chapter: completeness
% Section: modalities-ccs

\documentclass[../../../include/open-logic-section]{subfiles}

\begin{document}

\olfileid{nml}{com}{mod}

\olsection{المؤثرات الجهية والمجموعات المتسقة التامة}

\begin{explain}
قد اخترنا للنموذج~$\mModel{M^\Sigma}$ عوالم هي المجموعات
$\Sigma$-المتسقة التامة~$\Delta$ في منطق جهوي نظامي~$\Sigma$.
وبقي أن نعيّن علاقة إمكان الوصول~$R^\Sigma$ بين هذه «العوالم».
وضابط الاختيار، في هذه العلاقة وفي التعيين~$V^\Sigma$، أن يتحقق
$\mSat{M^\Sigma}{!A}[\Delta]$ إذا وفقط إذا كان~$!A \in \Delta$.
فلننظر كيف يهدي هذا الضابط إلى التعريف.

متى عُرّفت علاقة إمكان الوصول، ضمن تعريف الصدق عند عالم أن
\iftag{prvBox}{$\mSat{M^\Sigma}{\Box!A}[\Delta]$ إذا وفقط إذا كان
  $\mSat{M^\Sigma}{!A}[\Delta']$ لكل~$\Delta'$ بحيث
  $R^\Sigma\Delta\Delta'$}{$\mSat{M^\Sigma}{\Diamond!A}[\Delta]$ إذا
  وفقط إذا كان~$\mSat{M^\Sigma}{!A}[\Delta']$ لبعض~$\Delta'$ بحيث
  $R^\Sigma\Delta\Delta'$}. ويتطلب برهان أن
$\mSat{M^\Sigma}{!A}[\Delta]$ إذا وفقط إذا كان~$!A \in \Delta$ أن
يصدق ذلك خصوصًا في ال!!{formula}s التي تبدأ بمؤثر جهوي؛ أي إن
\iftag{prvBox}{$\mSat{M^\Sigma}{\Box!A}[\Delta]$ إذا وفقط إذا كان
  $\Box !A \in \Delta$}{$\mSat{M^\Sigma}{\Diamond!A}[\Delta]$ إذا
  وفقط إذا كان~$\Diamond !A \in \Delta$}. وبجمع هذا المطلب مع تعريف
الصدق عند عالم لـ\iftag{prvBox}{$\Box !A$}{$\Diamond !A$} نحصل على:
\[
\iftag{prvBox}{%
  \Box!A \in \Delta \text{ إذا وفقط إذا كان } !A \in \Delta'
  \text{ لكل $\Delta'$ يحقق } R^\Sigma\Delta\Delta'}{%
  \Diamond!A \in \Delta \text{ إذا وفقط إذا كان } !A \in \Delta'
  \text{ لبعض } \Delta' \text{ يحقق } R^\Sigma\Delta\Delta'}
\]
\iftag{prvBox}{تأمل الاتجاه من اليسار إلى اليمين: فهو يقول إنه إذا كان
  $\Box !A \in \Delta$، كان~$!A \in \Delta'$ لأي~$!A$ وأي~$\Delta'$
  يحقق~$R^\Sigma\Delta\Delta'$. فإذا اشترطنا أن
  $R^\Sigma\Delta\Delta'$ إذا وفقط إذا كان~$!A \in \Delta'$ لكل
  $\Box!A \in \Delta$، تحقق ذلك. ويمكننا أن نكتب الشرط الواقع يمين
  ``إذا وفقط إذا'' بإيجاز أكبر هكذا:
  $\Setabs{!A}{\Box!A \in \Delta} \subseteq \Delta'$.}{تأمل الاتجاه
  من اليمين إلى اليسار: فهو يقول إنه إذا كان~$!A \in \Delta'$، كان
  $\Diamond!A \in \Delta$ لأي~$!A$ وأي~$\Delta'$ يحقق
  $R^\Sigma\Delta\Delta'$.
  % تصحيح مصدر (0445-I1): التعريف المعلن علاقة تكافؤ لا شرط كفاية وحده.
  فإذا عرّفنا~$R^\Sigma\Delta\Delta'$
  بأنه يتحقق إذا وفقط إذا كان~$\Diamond!A \in \Delta$ لكل~$!A \in \Delta'$،
  لزم هذا الاتجاه من التعريف.
  ويمكننا أن نكتب الشرط الواقع يمين ``إذا وفقط إذا'' بإيجاز أكبر هكذا:
  $\Setabs{\Diamond!A}{!A \in \Delta'} \subseteq \Delta$.}

ويبقى السؤال: هل يضمن تعريف~$R^\Sigma$ هذا بالفعل أن
\iftag{prvBox}{$\Box !A \in \Delta$ إذا وفقط إذا كان
  $\mSat{M^\Sigma}{\Box !A}[\Delta]$؟ \iftag{prvDiamond}{وهل يضمن
    أيضًا أن }{}}{}\iftag{prvDiamond}{$\Diamond !A \in \Delta$ إذا
  وفقط إذا كان~$\mSat{M^\Sigma}{\Diamond !A}[\Delta]$؟}{}
ستثبت النتائج القليلة الآتية ذلك.
\end{explain}

\begin{defn}
  إذا كانت~$\Gamma$ مجموعة من ال!!{formula}s، فلتكن
  \begin{align*}
    \Box\Gamma & = \Setabs{\Box !B}{!B \in \Gamma}\\
    \Diamond\Gamma & = \Setabs{\Diamond !B}{!B \in \Gamma}\\
    \intertext{ولتكن}
    \Box^{-1}\Gamma & = \Setabs{!B}{\Box !B \in \Gamma}\\
    \Diamond^{-1}\Gamma & = \Setabs{!B}{\Diamond !B \in \Gamma}\\
  \end{align*}
\end{defn}

فالمقصود بـ~$\Box\Gamma$ أخذ~$\Gamma$ ووضع~$\Box$ أمام كل
!!{formula} من~$\Gamma$. وأما~$\Box^{-1}\Gamma$ فنأخذ لتكوينها
ال!!{formula}s التي تبدأ بـ~$\Box$ في~$\Gamma$، ثم نحذف~$\Box$ الأولى
من كل واحدة. وإنما أدخلنا هذا الاصطلاح اختصارًا للترميز،
لا لأن في هذا التعريف وحده نتيجة جديدة مهمة.

لاحظ أن~$\Box\Box^{-1}\Gamma \subseteq \Gamma$:
\[
\Box\Box^{-1}\Gamma = \Setabs{\Box!B}{\Box!B \in \Gamma}
\]
أي إنها ليست إلا مجموعة جميع ال!!{formula}s في~$\Gamma$ التي تبدأ
بـ~$\Box$.

\begin{lem}\ollabel{lem:box1}
  إذا كان~$\Gamma \Proves[\Sigma] !A$، فإن
  $\Box\Gamma \Proves[\Sigma] \Box !A$.
\end{lem}

\begin{proof}
  من~$\Gamma \Proves[\Sigma] !A$ نختار مقدمات متناهية
  $!B_1$، \dots،~$!B_k \in \Gamma$ تحقق
  $\Sigma \Proves !B_1 \lif (!B_2 \lif \cdots (!B_k \lif !A)\cdots)$.
  ولأن~$\Sigma$ نظامية، يصح استعمال قاعدة~\RK{}، فنحصل على
  $\Sigma \Proves \Box!B_1 \lif (\Box!B_2 \lif \cdots
  (\Box!B_k \lif \Box!A)\cdots)$.
  والمقدمات الجديدة $\Box!B_1$، \dots،~$\Box!B_k \in \Box\Gamma$
  بحكم التعريف. فهذا الاشتقاق بعينه شاهد على
  $\Box\Gamma \Proves[\Sigma] \Box !A$.
\end{proof}

\begin{lem}\ollabel{lem:box2}
  إذا كان~$\Box^{-1}\Gamma \Proves[\Sigma] !A$، فإن
  $\Gamma \Proves[\Sigma] \Box!A$.
\end{lem}

\begin{proof}
  لنفرض $\Box^{-1}\Gamma \Proves[\Sigma] !A$.
  بتطبيق \olref{lem:box1} يلزم
  $\Box\Box^{-1}\Gamma \Proves[\Sigma] \Box!A$.
  وقد علمنا أن $\Box\Box^{-1}\Gamma \subseteq \Gamma$؛
  فيجوز توسيع المقدمات، وبالرتابة نحصل على
  $\Gamma \Proves[\Sigma] \Box!A$.
\end{proof}

\iftag{prvBox}{%
% Box is primitive

\begin{prop}\ollabel{prop:box}
  إذا كانت~$\Gamma$ $\Sigma$-متسقة تامة، فإن~$\Box !A \in \Gamma$
  إذا وفقط إذا كان، لكل مجموعة $\Sigma$-متسقة تامة~$\Delta$ تحقق
  $\Box^{-1}\Gamma \subseteq \Delta$، لدينا~$!A \in \Delta$.
\end{prop}

\begin{proof}
  افترض أن~$\Gamma$ $\Sigma$-متسقة تامة. الاتجاه ``فقط إذا'' يسير:
  افترض أن~$\Box !A \in \Gamma$ وأن~$\Box^{-1}\Gamma \subseteq
  \Delta$. لما كان~$\Box!A \in \Gamma$، كان
  $!A \in \Box^{-1}\Gamma \subseteq \Delta$، ومن ثم~$!A \in \Delta$.

  وأما اتجاه «إذا» فنأخذ مقابله. لنفرض $\Box!A \notin \Gamma$.
  بما أن~$\Gamma$ $\Sigma$-متسقة تامة، فهي مغلقة استنتاجيًا؛
  ومن ثم $\Gamma \Proves/[\Sigma] \Box !A$.
  ويقتضي مقابل~\olref{lem:box2} أن
  $\Box^{-1}\Gamma \Proves/[\Sigma] !A$.
  وبحسب \olref[prf][con]{prop:consistencyfacts}\olref[prf][con]{prop:consistencyfacts-b}
  تكون~$\Box^{-1}\Gamma \cup \{\lnot!A\}$ $\Sigma$-متسقة.
  فنوسّعها بمبرهنة ليندنباوم إلى مجموعة $\Sigma$-متسقة تامة~$\Delta$
  تحقق $\Box^{-1}\Gamma \cup \{\lnot!A\} \subseteq \Delta$.
  ولما كانت مشتملة على النفي ومتسقة، كان $!A \notin \Delta$،
  فوجدنا المجموعة التي تنقض الشرط الكلي.
\end{proof}
}{%
% Box is not primitive, only Diamond is

\begin{prop}\ollabel{prop:diamond}
  إذا كانت~$\Gamma$ $\Sigma$-متسقة تامة، فإن~$\Diamond !A \in \Gamma$
  إذا وفقط إذا وجدت مجموعة $\Sigma$-متسقة تامة~$\Delta$ تحقق
  $\Diamond\Delta \subseteq \Gamma$ و~$!A \in \Delta$.
\end{prop}

\begin{proof}
  افترض أن~$\Gamma$ $\Sigma$-متسقة تامة. الاتجاه من اليمين إلى اليسار
  يسير: لتكن~$\Delta$ مجموعة $\Sigma$-متسقة تامة تحقق
  $\Diamond\Delta \subseteq \Gamma$ و~$!A \in \Delta$. عندئذ
  $\Diamond!A \in \Diamond\Delta \subseteq \Gamma$، أي
  $\Diamond!A \in \Gamma$.

  أما الاتجاه من اليسار إلى اليمين، فافترض أن
  $\Diamond !A \in \Gamma$. علينا أن نبين وجود مجموعة $\Sigma$-متسقة
  تامة~$\Delta$ بحيث~$!A \in \Delta$ و
  $\Diamond\Delta \subseteq \Gamma$. لما كان~$\Diamond !A \in \Gamma$
  وكانت~$\Gamma$ $\Sigma$-متسقة، كان~$\Box\lnot !A \notin \Gamma$.
  ولما كانت~$\Gamma$ $\Sigma$-متسقة تامة، كانت مغلقة استنتاجيًا، ومن
  ثم~$\Gamma \Proves/[\Sigma] \Box\lnot !A$. وبحسب~\olref{lem:box2}،
  $\Box^{-1}\Gamma \Proves/[\Sigma] \lnot !A$. وبحسب
  \olref[prf][con]{prop:consistencyfacts}\olref[prf][con]{prop:consistencyfacts-b}،
  تكون~$\Box^{-1}\Gamma \cup \{!A\}$ $\Sigma$-متسقة. وبمبرهنة
  ليندنباوم، توجد مجموعة $\Sigma$-متسقة تامة~$\Delta$ بحيث
  $\Box^{-1}\Gamma \cup \{!A\} \subseteq \Delta$. ومن الواضح أن
  $!A \in \Delta$. ولإثبات~$\Diamond\Delta \subseteq \Gamma$، افترض
  خلافه: يوجد~$!B \in \Delta$ بحيث~$\Diamond !B \notin \Gamma$.
  ولما كانت~$\Gamma$ $\Sigma$-متسقة تامة، كان
  $\Box\lnot !B \in \Gamma$. لكن عندئذ أيضًا
  $\lnot !B \in \Box^{-1}\Gamma \subseteq \Delta$، وهذا يجعل
  $\Delta$ غير متسقة.
\end{proof}
}

\begin{lem}\ollabel{lem:box-iff-diamond}
  افترض أن~$\Gamma$ و$\Delta$ مجموعتان $\Sigma$-متسقتان تامتان. عندئذ
  $\Box^{-1}\Gamma \subseteq \Delta$ إذا وفقط إذا كان
  $\Diamond\Delta \subseteq \Gamma$.
\end{lem}

\begin{proof}
  لإثبات اتجاه «فقط إذا»، نفرض~$\Box^{-1}\Gamma \subseteq \Delta$،
  ونأخذ~$\Diamond!A \in \Diamond\Delta$، أي~$!A \in \Delta$.
  المطلوب $\Diamond!A \in \Gamma$؛ ويكفي له~$\Box\lnot!A \notin \Gamma$،
  إذ تلزم عندئذ~$\lnot\Box\lnot !A \in \Gamma$ بالقصوية.
  ولو كان~$\Box\lnot!A \in \Gamma$ لكان
  $\lnot!A \in \Delta$ بفرض الاحتواء، فيبطل اتساق~$\Delta$،
  لأن~$!A \in \Delta$. فثبت~$\Box\lnot!A \notin \Gamma$ والمطلوب.

  ولإثبات اتجاه «إذا»، نفرض~$\Diamond\Delta \subseteq \Gamma$
  ونستعمل المقابل: نأخذ~$!A \notin \Delta$ ونثبت~$\Box!A \notin \Gamma$.
  فمن~$!A \notin \Delta$ يلزم~$\lnot!A \in \Delta$ بالقصوية؛
  وبفرض الاحتواء يلزم~$\Diamond\lnot!A \in \Gamma$.
  غير أن~$\Diamond\lnot!A$ تكافئ~$\lnot\Box !A$ في كل منطق جهوي نظامي.
  فتكون هذه الأخيرة في~$\Gamma$ بالإغلاق، ويمنع الاتساق خلافها:
  $\Box!A \notin\Gamma$.
\end{proof}

\iftag{prvBox}{%
  \iftag{prvDiamond}{%
% Box and Diamond are both primitive; prove
% prop:diamond using lem:box-iff-diamond

\begin{prop}\ollabel{prop:diamond}
  إذا كانت~$\Gamma$ $\Sigma$-متسقة تامة، فإن~$\Diamond !A \in \Gamma$
  إذا وفقط إذا وجدت مجموعة $\Sigma$-متسقة تامة~$\Delta$ تحقق
  $\Diamond\Delta \subseteq \Gamma$ و~$!A \in \Delta$.
\end{prop}

\begin{proof}
افترض أن~$\Gamma$ $\Sigma$-متسقة تامة. لدينا
$\Diamond!A \in \Gamma$ إذا وفقط إذا كان
$\lnot\Box\lnot !A \in \Gamma$، بحسب~\Dual{} والإغلاق.
ولدينا~$\lnot\Box\lnot !A \in \Gamma$ إذا وفقط إذا كان
$\Box\lnot !A \notin \Gamma$، بحسب
\olref[ccs]{prop:ccs-properties}\olref[ccs]{prop:ccs-lnot}، لأن
$\Gamma$ $\Sigma$-متسقة تامة. وبحسب~\olref{prop:box}، يكون
$\Box\lnot !A \notin \Gamma$ إذا وفقط إذا وجدت مجموعة $\Sigma$-متسقة
تامة~$\Delta$ تحقق~$\Box^{-1}\Gamma \subseteq \Delta$ و
$\lnot!A \notin \Delta$. تأمل الآن أي~$\Delta$ كهذه. بحسب
\olref{lem:box-iff-diamond}، يكون~$\Box^{-1}\Gamma \subseteq \Delta$
إذا وفقط إذا كان~$\Diamond\Delta \subseteq \Gamma$. كذلك، يكون
$\lnot !A \notin \Delta$ إذا وفقط إذا كان~$!A \in \Delta$، بحسب
\olref[ccs]{prop:ccs-properties}\olref[ccs]{prop:ccs-lnot}. وهكذا،
$\Diamond!A \in \Gamma$ إذا وفقط إذا وجدت مجموعة $\Sigma$-متسقة
تامة~$\Delta$ تحقق~$\Diamond\Delta \subseteq \Gamma$ و
$!A \in \Delta$.
\end{proof}

}{}}{}

\begin{prob}
  بيّن أنه إذا كانت~$\Gamma$ $\Sigma$-متسقة تامة، فإن
  \iftag{prvBox}{$\Diamond !A \in \Gamma$ إذا وفقط إذا وجدت مجموعة
    $\Sigma$-متسقة تامة~$\Delta$ بحيث
    $\Box^{-1}\Gamma \subseteq \Delta$ و~$!A \in \Delta$.}{$\Box !A
    \in \Gamma$ إذا وفقط إذا كان، لكل مجموعة $\Sigma$-متسقة تامة
    $\Delta$ تحقق~$\Box^{-1}\Gamma \subseteq \Delta$، لدينا
    $!A \in \Delta$.} \emph{افعل ذلك من دون استعمال
    \olref[nml][com][mod]{lem:box-iff-diamond}}.
\end{prob}

\end{document}
